% FORMALIZACION_REUNIDA: representación pentacomponente de II,
% con antecedente Paley y demostración directa de los momentos del marco.
% El retorno transporta la carta; no se identifica un ciclo de orden doce con A5.
\section{Incidence orientée et représentation tensorielle gravitationnelle}
\label{vii:sec:pentacomponente}
La représentation exceptionnelle de l'état conserve l'incidence qui accompagne sa lecture arithmétique. La matrice ternaire \(A_W\) de \eqref{exc:aw} fournit ici un antécédent matériel précis. Son relèvement équilibré \(C=\operatorname{bal}(A_W)\), avec les représentants \(0,1,-1\), satisfait
\[
 C=C^{\mathsf T},\qquad C^2=5I_6,\qquad \operatorname{Tr}C=0.
\]
Les six positions ont pour ordre \((\infty,0,1,2,3,4)\). Cette incidence produit d'abord un espace tridimensionnel ; sa lecture quadratique produit ensuite l'espace à cinq composantes sur lequel agit le couplage gravitationnel.

\subsection{Revêtement signé et transport du calendrier}
On définit
\[
 E_C=\frac12(I_6+C/\sqrt5),\quad
 \mathcal H_C=\operatorname{im}E_C,\quad
 v_{u,\sigma}=\sigma\sqrt2E_Ce_u,\quad \sigma\in\{1,-1\}.
\]
\begin{proposition}[Douze positions orientées]
\label{vii:prop:doce-orientadas}
Les vecteurs \(v_{u,\sigma}\) sont douze vecteurs unitaires distincts
dans un espace de dimension trois. Leur incidence est déterminée par
\(\langle v_{u,\sigma},v_{v,\tau}\rangle=1/\sqrt5\).
L'inversion de l'orientation agit par antipodie.
\end{proposition}
\begin{proof}
L'identité quadratique de \(C\) implique \(E_C^2=E_C=E_C^*\) ; sa trace est trois. Puisque la diagonale de \(C\) est nulle,
\[
 \langle v_{u,\sigma},v_{v,\tau}\rangle
 =\begin{cases}\sigma\tau,&u=v,\\
 \sigma\tau C_{uv}/\sqrt5,&u\ne v.
 \end{cases}
\]
Cela prouve la normalisation, la distinction des douze vecteurs et l'incidence.
Remplacer \(\sigma\) par \(-\sigma\) change le signe du vecteur.
\end{proof}

Une expression euclidienne explicite s'obtient en posant
\(\nu=\sqrt{1+\varphi^2}\) et
\[
 B=\frac1\nu
 \begin{pmatrix}
 0&-1&-\varphi&0&\varphi&1\\
 -1&-\varphi&0&1&0&-\varphi\\
 -\varphi&0&-1&-\varphi&-1&0
 \end{pmatrix}.
\]
La relation \(\varphi^2=\varphi+1\), pour la coordonnée déjà générée, donne \(B^*B=2E_C\). En conséquence, \(B/\sqrt2|_{\mathcal H_C}\) est une isométrie et envoie \(v_{u,\sigma}\) sur \(\sigma Be_u\). L'image est
\[
 \mathcal V=\frac1\nu\bigl(
 \{0\}\times\{\pm1\}\times\{\pm\varphi\}\ \cup\
 \{\pm1\}\times\{\pm\varphi\}\times\{0\}\ \cup\
 \{\pm\varphi\}\times\{0\}\times\{\pm1\}\bigr).
\]
La coordonnée interne, l'incidence et leurs
douze représentants géométriques sont ainsi reliés.

La coordonnée du calendrier s'écrit \((b,r)\in
\mathbb Z/12\mathbb Z\times\{0,\ldots,8\}\). Une route de longueur \(n\)
induit
\[
 k(r,n)=\left\lfloor\frac{r+n}{9}\right\rfloor,\qquad
 (b,r)\longmapsto
 \bigl(b+k(r,n)\bmod12,\ (r+n)\bmod9\bigr).
\]
L'alignement explicite de ses positions \(p=b+1\) avec la coordinatisation signée est
\[
\begin{array}{c|rrrrrrrrrrrr}
p&1&2&3&4&5&6&7&8&9&10&11&12\\ \hline
u&\infty&0&0&\infty&1&3&3&1&2&4&4&2\\
\sigma&+&-&+&-&-&+&-&+&-&+&-&+
\end{array}
\]
Ce tableau fixe une trivialisation du porteur. Chaque changement de position
transporte conjointement les vecteurs et leur matrice d'incidence.

\begin{proposition}[Covariance du repère sous retour]
Si \(S e_p=e_{p+1\bmod12}\), la représentation matricielle d'un opérateur \(P\) du support se transporte comme \(P_k=S^kPS^{-k}\). Ce transport conserve les produits scalaires, les incidences et la composition des parcours. Après douze retours, la coordinatisation est rétablie, tandis que la mémoire accumulée demeure conservée dans l'état source.
\end{proposition}
\begin{proof}
\(S\) est orthogonale et \(S^{12}=I\) ; les propriétés métriques et la périodicité de la coordinatisation en découlent. Pour la composition, avec \(r'=(r+n)\bmod9\), la division euclidienne donne \(k(r,n+m)=k(r,n)+k(r',m)\). Les exposants de \(S\) s'additionnent exactement avec la retenue. L'opération agit sur la coordonnée du calendrier ; les mises à jour de mémoire du parcours restent présentes dans son domaine enrichi.
\end{proof}

\subsection{Lecture quadratique et espace à cinq composantes}
Soient \(\mathbb R^{\mathcal V}\) l'espace des douze coordonnées,
\(\mathsf A e_v=e_{-v}\) son involution antipodale et
\(\mathsf J=\mathbf1\mathbf1^{\mathsf T}\). On définit
\[
 P_5=\frac12(I+\mathsf A)-\frac1{12}\mathsf J,\qquad
 Q_v=vv^{\mathsf T}-\frac13I_3,\qquad
 \Gamma_0x=\sum_{v\in\mathcal V}x_vQ_v.
\]
\begin{theorem}[Projecteur à cinq composantes et isométrie tensorielle]
\label{vii:thm:portador-pentacomponente}
\(P_5\) est le projecteur orthogonal de rang cinq sur l'espace des
coordonnées paires sous l'antipodie et de somme nulle. De plus,
\[
 \Gamma_0^*\Gamma_0=\frac85P_5,\qquad
 \Gamma_0\Gamma_0^*=\frac85I_{\operatorname{Sym}^2_0(\mathbb R^3)}.
\]
Par conséquent
\[
 \Gamma_5=\sqrt{\frac58}\,\Gamma_0|_{\operatorname{im}P_5}
\]
est une isométrie sur l'espace des tenseurs symétriques sans trace.
\end{theorem}
\begin{proof}
L'espace pair possède une coordonnée pour chaque paire antipodale, donc
sa dimension est six. Le vecteur constant lui appartient ; soustraire son
projecteur de rang un donne \(P_5\), de rang cinq.

Comme
\(\langle Q_v,Q_w\rangle_F=(v^{\mathsf T}w)^2-1/3\), les entrées
de \(\Gamma_0^*\Gamma_0\) valent \(2/3\) pour \(w=\pm v\) et
\(-2/15\) aux dix autres emplacements de chaque ligne. Ce sont
exactement les entrées de \((8/5)P_5\).

On peut également vérifier l'identité dans l'espace tensoriel sans recourir à une classification des représentations. La liste des sommets donne
\[
 \sum_vv_i^4=\frac{12}{5},\qquad
 \sum_vv_i^2v_j^2=\frac45\quad(i\ne j),
\]
et les moments comportant une puissance impaire s'annulent par les changements de signe. En effet, \((1+\varphi^2)^2=5\varphi^2\) et \(1+\varphi^4=3\varphi^2\). Pour un tenseur symétrique \(Q\), ces sommes donnent
\[
 \sum_v(v^{\mathsf T}Qv)\,vv^{\mathsf T}
 =\frac85Q+\frac45(\operatorname{tr}Q)I.
\]
Avec \(\operatorname{tr}Q=0\), et en utilisant
\(\sum_vv^{\mathsf T}Qv=4\operatorname{tr}Q=0\), on obtient
\(\Gamma_0\Gamma_0^*Q=(8/5)Q\). Les deux identités prouvent
l'isométrie normalisée et sa surjectivité.
\end{proof}

La section gravitationnelle de \(\ref{v:ant:sec:gravity}\) détermine
l’opérateur \(G_aP_5\). Son spectre a \(G_a\) de multiplicité
cinq et zéro de multiplicité sept. L’orientation est conservée dans
le porteur et dans son transport ; \(G_a\) fixe le couplage commun.

\subsection{Projection transversale et hélicités}
Pour \(n\in S^2\), posons \(\Pi=I-nn^{\mathsf T}\) et
\[
 \Lambda_nQ=\Pi Q\Pi-\frac12\operatorname{tr}(\Pi Q\Pi)\Pi.
\]
\begin{proposition}[Réduction transverse de rang deux]
\(\Lambda_n\) est le projecteur orthogonal sur \(\{Q\in\operatorname{Sym}^2_0(\mathbb R^3):Qn=0\}\), de dimension deux.
\end{proposition}
\begin{proof}
\(\Pi n=0\) donne la transversalité et \(\operatorname{tr}\Pi=2\) donne
une trace nulle. Un tenseur transverse sans trace reste fixe.
La formule est autoadjointe pour le produit de Frobenius.
Dans un repère \((e_1,e_2,n)\), l'image a pour base orthonormée
\[
 E_+=\frac{e_1e_1^{\mathsf T}-e_2e_2^{\mathsf T}}{\sqrt2},\qquad
 E_\times=\frac{e_1e_2^{\mathsf T}+e_2e_1^{\mathsf T}}{\sqrt2}.
\]
\end{proof}
Les trois autres vecteurs d'une base orthonormée de l'espace porteur sont
\[
 E_1=\frac{e_1n^{\mathsf T}+ne_1^{\mathsf T}}{\sqrt2},\quad
 E_2=\frac{e_2n^{\mathsf T}+ne_2^{\mathsf T}}{\sqrt2},\quad
 E_0=\frac{2nn^{\mathsf T}-e_1e_1^{\mathsf T}-e_2e_2^{\mathsf T}}{\sqrt6}.
\]
Lorsque le repère transversal tourne d'un angle \(\theta\), les paires \((E_+,E_\times)\) et \((E_1,E_2)\) tournent de \(2\theta\) et \(\theta\), respectivement, tandis que \(E_0\) reste fixe. La complexification possède ainsi les poids \(2,-2,1,-1,0\). \(\Lambda_n\) conserve les deux premiers et annule les autres. On établit ainsi la chaîne
\[
 \mathbb R^{12}\xrightarrow{P_5}\operatorname{im}P_5
 \xrightarrow{\Gamma_5}\operatorname{Sym}^2_0(\mathbb R^3)
 \xrightarrow{\Lambda_n}\mathcal H_{\rm TT}(n),
 \qquad 12\longrightarrow5\longrightarrow5\longrightarrow2.
\]
La représentation à cinq composantes et sa réduction sont des résultats
algébriques antérieurs au choix des équations de champ.
Dans la réalisation gravitationnelle sans masse, les contraintes de jauge
sélectionnent l'image transverse. Cette distinction conserve tant les
cinq coordonnées structurelles que les deux hélicités radiatives.
